\section{Local separation at distinct ordinates}\label{sec:localwindows}

The leading oscillation in \eqref{eq:gapsig} is small. Its sign nevertheless
changes frequently. A Gaussian turns that sign information into a lower bound
on the number of \emph{distinct} ordinates. For multisets $X,Y$ write
\[
 d_+(I)=\#\{\gamma\in I:m_X(\gamma)>m_Y(\gamma)\},\qquad
 d_-(I)=\#\{\gamma\in I:m_X(\gamma)<m_Y(\gamma)\}.
\]
Thus a shared ordinate can contribute one discrepancy if its multiplicities
differ. Throughout this section GRH for the input characters is assumed.

\begin{theorem}[Distinct discrepancies in short windows]\label{thm:localcoupled}
Compare two models of the same size. Suppose either their character multisets
differ, or their aligned labels are distinct and their principal minors differ.
There exist $c_0,C,T_0>0$ depending on the models such that for every
$T\ge T_0$ and $C\sqrt{\log\log T}\le H\le T/2$,
\[
 \min\{d_+([T-H,T+H]),d_-([T-H,T+H])\}\ \ge\ c_0 H .
\]
The interval even contains a sequence of $\gg H$ distinct ordinates, in increasing order, on which the direction of excess alternates.
For every $b_0>1/\pi$ each direction occurs at least once in
$[T-b_0\log\log\log T,T+b_0\log\log\log T]$ for all sufficiently large $T$.
\end{theorem}

\begin{proof}
Let $\nu=Z_{\cX,A}-Z_{\cY,B}$ as a signed point measure, and choose
$\sigma>1$, $b=\sigma-1/2$, $K_b(t)=b/(b^2+t^2)$. Write $H_b=K_b*\nu$.
The Hadamard identity \eqref{poiss} is valid on this line.
The calculation of \eqref{eq:gapsig} also holds for any fixed $\sigma>1$:
all Euler series and their derivatives converge absolutely there. In the
same-label case it gives
\begin{equation}\label{loc-asymp}
 H_b(t)=(\log t)^{-j}f_\sigma(t)+O((\log t)^{-j-1}),\qquad t\to+\infty ,
\end{equation}
with $f_\sigma$ real uniformly almost periodic, nonzero, and of mean zero.
The nonzero Dirichlet coefficient used there remains nonzero on every such
line. For unequal character multisets set $j=0$ and subtract \eqref{cline}
on the line $\Re s=\sigma$. Now
\[
 f_\sigma(t)=\tfrac12\log(Q_{\cX}/Q_{\cY})
 +\Re\!\sum_{m\ge2}\Lambda(m)(S_{\cY}(m)-S_{\cX}(m))m^{-\sigma-it}.
\]
In both cases $H_b(t)=O(\log(2+|t|))$ globally.

We first transfer \eqref{loc-asymp} to a positive test $w$. Use the transform
$\widehat w(\xi)=\int w(t)e^{-i\xi t}\,dt$. We shall use two even,
unit-mass densities: a Gaussian $w(t)=(\lambda/\pi)^{1/2}e^{-\lambda t^2}$,
and later a normalized $e^{-L\cosh(kt)}$. For the latter take
$k<\pi/(2b)$. For either choice, define an integrable real kernel $h$ by
\[
 \widehat h(\xi)=\pi^{-1}e^{b|\xi|}\widehat w(\xi).
\]
Indeed for the Gaussian its transform and all derivatives have Gaussian decay.
In the second case choose $b<\beta<\pi/(2k)$ and shift the Fourier contour for
$w(t)$, $tw(t)$ and $t^2w(t)$ to imaginary height $\pm\beta$. There their
absolute values decay doubly exponentially since $\cos(k\beta)>0$.
Thus the multiplier and its first two derivatives are exponentially bounded
at both infinities. It is continuous and piecewise smooth at zero.
Twice integrating its inverse Fourier integral by parts on the two half-lines
gives $h(v)=O((1+|v|)^{-2})$, allowing the jump of the first derivative.
As $\widehat K_b=\pi e^{-b|\xi|}$, we have $h*K_b=w$.

The unit-interval zero bound permits absolute convolution with $|h|*K_b$:
$K_b*|\nu|(u)\ll\log(2+|u|)$ by \eqref{poiss} for the two positive measures.
Consequently $w*\nu=h*H_b$. In the integral over $|v|\le\sqrt{t}$, replace
$\log(t-v)$ by $\log t+O(t^{-1/2})$; the discarded tails are
$O(t^{-1/2}\log(2+t))$ using the bound on $h$ and the global logarithmic
growth. We obtain
\begin{equation}\label{loc-smoothed}
 (w*\nu)(t)=(\log t)^{-j}G_w(t)+O((\log t)^{-j-1}),\qquad
 G_w=h*f_\sigma .
\end{equation}
The function $G_w$ is real uniformly almost periodic. For common labels its
mean is zero. The Gaussian multiplier never vanishes, so $G_w$ has a nonzero
Bohr coefficient. For the second family the normalized densities concentrate
at zero as $L\to\infty$, whence $\widehat w(\xi)\to1$ for each fixed $\xi$.
Taking $L$ sufficiently large preserves any one nonzero coefficient.

For different label multisets there can be a constant term. We verify both
signs explicitly. Let $\delta(p)=S_{\cY}(p)-S_{\cX}(p)$.
The mean of $G_w$ is $\log(Q_{\cX}/Q_{\cY})/(2\pi)$, and its coefficient
at $-\log p$ is $\delta(p)(\log p)\widehat w(\log p)/(2\pi\sqrt p)$.
The other nonconstant frequencies are signed logarithms of prime powers.
By the Riesz product in the proof of Theorem~\ref{main}, the means of $G_w$
times two nonnegative trigonometric polynomials (phases chosen in the two
opposite directions) are
\[
 \frac1{2\pi}\left(\log(Q_{\cX}/Q_{\cY})\ \pm
 \sum_{p\in S}\frac{\log p}{\sqrt p}
       |\delta(p)\widehat w(\log p)|\right).
\]
To check the factor, the coefficients at frequencies $\pm\log p$ in a Riesz
product each have modulus $1/2$ and the two conjugate terms add.
Choose the finite prime set to exceed the absolute conductor term by a
strict margin when $\widehat w=1$. Then choose the Gaussian parameter
$\lambda$ (or the parameter $L$) sufficiently large to preserve the margin.
This gives a positive and a negative value of $G_w$. In the common-label
case the same conclusion follows from nonzero, zero-mean almost periodicity:
a nonzero nonnegative uniformly almost-periodic function has positive mean
by uniform continuity and relatively dense approximate periods.

There are therefore numbers $a_0>0$ and $D<\infty$ such that every interval
of length $D$ contains points with $G_w>2a_0$ and with $G_w<-2a_0$.
Indeed, translate one point of each sign by relatively dense uniform
approximate periods. In successive blocks of length $D$ we may select samples
of alternating sign. Their values in \eqref{loc-smoothed} have magnitude at
least $a_0(\log t)^{-j}$ at large heights.

Take first the Gaussian and $I=[T-H,T+H]$. At $t$ in its middle half the
contribution to $w*\nu$ from $\gamma\notin I$ is
$O(\log(2+T)e^{-c_1H^2})$ for some $c_1>0$.
This follows by the unit-interval bound, summing
$(\log(2+T)+\log(2+m))e^{-\lambda m^2}$ at distances
$m\ge H/2-1$. For $H\ge C\sqrt{\log\log T}$ with $C$ large, the error is
$o((\log T)^{-j})$. Hence the Gaussian sum with just the atoms in $I$ has
$\gg H$ alternating samples and thus $\gg H$ distinct real zeros between them.

Here is the elementary form of Gaussian variation diminution that we need.
Multiplying this finite sum by the positive $e^{\lambda t^2}$ gives a
constant times $E(t)=\sum_\gamma b_\gamma e^{2\lambda\gamma t}$ with
$\operatorname{sign} b_\gamma=\operatorname{sign}\nu(\{\gamma\})$ at
each discrepancy. In an exponential sum with increasing real exponents,
the number of distinct real zeros is at most the number of sign changes of
the ordered nonzero coefficients. To prove this by induction, the case of
no sign change is immediate. Otherwise take a real $q$ between two
successive exponents at a change of sign. The operation
$E\mapsto e^{qt}(e^{-qt}E)'$ multiplies each coefficient by its exponent
minus $q$, reversing the entire initial block and removing exactly one
sign change. Rolle's theorem gives at least one fewer zero than before,
which proves the claim. Applied here it forces $\gg H$ sign changes among
the ordered atoms of $I$, hence that many alternating discrepancies and
the asserted bound in both directions. This finite argument is an
instance of \href{https://doi.org/10.1073/pnas.34.4.164}{Schoenberg's
variation-diminishing property}.

For the shorter window choose $\sigma>1$ sufficiently close to $1$ and
$k<\pi/(2b)$ sufficiently close to its upper endpoint that $kb_0>1$.
Use the second positive density above. The positive and negative sample
points can both be taken within $D$ of $T$. For either point the mass
contributed by ordinates outside the proposed window is
\[
 O\left(\log(2+T)\exp\{-c_2\exp(k b_0\log\log\log T)\}\right)
       =o((\log T)^{-j}).
\]
Here $D$ is absorbed in $c_2>0$ and the same unit-interval sum is used.
The signs therefore persist after truncation. Since $w$ is positive, there
must be an atom of each sign inside this window. This proves the theorem.
\end{proof}

\begin{corollary}[Blind recovery from positive ordinates]\label{cor:blindall}
Let the input list consist of $n$ distinct primitive characters.
Allow arbitrary parities and conductors, and any competing list of size at most
$n$. If
$\liminf_{T\to\infty}d_+((0,T])/T=0$, then the lists coincide up to
permutation and their corresponding principal minors agree. For weighted
graph Laplacians with nonnegative weights this identifies every edge and weight.
\end{corollary}
\begin{proof}
For fewer competing labels, \eqref{count} gives $\gg T\log T$ missing multiplicities. The unit-interval bound gives at most $O(\log(2+T))$ at each ordinate, forcing $d_+((0,T])\gg T$. At equal size, if the multisets of labels differ, or if an aligned minor differs after the labels agree, apply Theorem~\ref{thm:localcoupled} with center
$2T/3$ and radius $T/3$. Its interval lies in $(0,T]$ and contains
$\gg T$ distinct positive discrepancies. This contradicts the hypothesis.
Finally the one- and two-vertex minors determine $A_{ii}$ and $|A_{ij}|$,
and the off-diagonal signs of a graph Laplacian are fixed.
\end{proof}

The article obtains triple-logarithmic windows for zeros of products of
ordinary Dirichlet functions (\texttt{rec:thm:bothsigns}). Its result for
two distinct differential operators on one completed input has radius
$O(\log\log T)$ (\texttt{rec:cor:operatorlocal}). Here the unknown is a
Hermitian determinant coupling; the informative prime oscillation can first
appear in a negative power of $\log T$. Gaussian variation diminution
converts these shrinking oscillations into a linear number of distinct
multiplicity discrepancies.
